%% ─────────────────────────────────────────────────────────────────────────────
%% Capítulo 7 — Conclusiones y Trabajo Futuro
%% ─────────────────────────────────────────────────────────────────────────────

\chapter{Conclusiones y Trabajo Futuro}
\label{cap:conclusiones}

% Párrafo de apertura: síntesis global del trabajo realizado en 1-2 párrafos.
% ¿Se han cumplido los objetivos del Cap. 1? ¿Cuál es la aportación principal?

\todo{Escribir párrafo de apertura del capítulo.}

% ─────────────────────────────────────────────────────────────────────────────
\section{Conclusiones}
\label{sec:conclusiones}

% Guía:
% Estructurar las conclusiones alrededor de los 5 objetivos específicos
% del Cap. 1 (Sección~\ref{sec:objetivos}).
% Para cada objetivo: ¿se ha alcanzado? ¿Cuáles son los resultados clave?

\todo{Escribir sección de conclusiones.}

% ─────────────────────────────────────────────────────────────────────────────
\section{Limitaciones}
\label{sec:limitaciones}

% Guía:
% 1. pit_loss y P(SC) aún ruidosos; fallback hardcoded reduce calidad.
% 2. Solo temporada 2023 (OpenF1 desde 2023).
% 3. No incorpora dinámica de tráfico más allá del ruido gaussiano.
% 4. No modela estrategias de undercut/overcut activo.
% 5. No considera el impacto de la carga de combustible.
% 6. La pestaña Live de la web no está implementada.
% 7. I/O bloqueante en ingestión.
% 8. Los coeficientes OLS del perfil paramétrico (§3) y los rollouts del
%    Transformer (§2) operan sin garantías físicas en pilotos con muestra
%    reducida. La envoltura de validez (§4.4.5) los acota dentro del rango
%    histórico del circuito, lo que reduce la dispersión cross-driver de
%    8\,343 s a 375 s en Sakhir 2023, pero no resuelve el sobreajuste de
%    raíz; un perfil entrenado sobre datos limpios seguiría siendo
%    preferible a depender del clamp final.

\todo{Escribir sección de limitaciones.}

% ─────────────────────────────────────────────────────────────────────────────
\section{Trabajo futuro}
\label{sec:trabajo_futuro}

\subsection{Mejoras a corto plazo}
\label{ssec:mejoras_corto}

% - Migrar el cliente HTTP a httpx.AsyncClient para I/O no bloqueante.
% - Derivar pit_loss de los tiempos de parada en boxes (pit_out_time).
% - Derivar P(SC) de counts históricos de mensajes de race control.
% - Implementar caché TTLCache de cachetools en lugar del dict sin límite.
% - Añadir logging estructurado (structlog, JSON lines).

\todo{Escribir subsección de mejoras a corto plazo.}

\subsection{Extensiones del modelo}
\label{ssec:extensiones_modelo}

% - Incorporar más temporadas (2024, 2025) cuando OpenF1 las incluya.
% - Explorar variantes ligeras del Transformer (Linformer, Performer) para reducir coste.
% - Explorar N-BEATS o PatchTST como baseline.
% - Integrar modelos fundacionales como TimesFM.
% - Modelar el undercut/overcut con un modelo de tráfico explícito.
% - Añadir cuantificación de incertidumbre mediante dropout bayesiano.

\todo{Escribir subsección de extensiones del modelo.}

\subsection{Evolución del sistema}
\label{ssec:evolucion_sistema}

% - Implementar las pestañas Live y Rewatch del frontend.
% - Modelar interacciones estratégicas entre pilotos.
% - Multi-agent RL para estrategias competitivas.
% - Despliegue en cloud (Docker + Kubernetes).
% - Integrar datos de otras categorías (F2, Formula-E) para transfer learning.

\todo{Escribir subsección de evolución del sistema.}
